\chapter{Core Experiments: The Canonical Setup}\label{sec:experiments}

\section{Experiments summary}\label{sec:experiments-summary}

Two experiments constitute the core of this chapter: the canonical relay comparison (E2) and the parameter-normalization and complexity study (E3), which together answer every canonical hypothesis (H1--H4). A supplementary coded study (Section~\ref{sec:coded-block-df}) departs from the canonical setup within this same chapter, adding a rate-1/2 convolutional code, a block-DF relay and an adaptive modulation-and-coding scheme. The principal contribution (Chapter~\ref{sec:unknown-channels}) is the unknown-channel study (E6), which tests H5.
The core (E2, E3) and unknown-channel (E6) experimental configurations are summarised in the table below; the supplementary coded study's configuration is described separately in Section~\ref{sec:coded-block-df}.

Each experiment subsection reports the key results inline (BER tables and figures); supporting material and model hyperparameters are collected in the appendices (Appendices~\ref{sec:appendix-b-model-architectures-and-hyperparameters}--\ref{sec:appendix-c-software-architecture}).

\begin{longtable}{p{0.9cm} p{2.2cm} p{2.1cm} p{3.6cm} p{3.9cm}}
\caption[Summary of experimental configurations]{Summary of core and unknown-channel experimental configurations}
\label{tab:experiment-summary}\\
\hline
\textbf{Exp.} &
\textbf{Topology} &
\textbf{Modulation} &
\textbf{Channel / Equalization} &
\textbf{Focus} \\
\hline
\endfirsthead

\hline
\textbf{Exp.} &
\textbf{Topology} &
\textbf{Modulation} &
\textbf{Channel / Equalization} &
\textbf{Focus} \\
\hline
\endhead

E2 & SISO & QPSK &
Rayleigh (canonical) &
Canonical relay comparison (H1, H2) \\
\hline

E3 & SISO & QPSK &
Rayleigh (canonical) &
Parameter normalization \& complexity (H3, H4) \\
\hline

E6 & SISO & BPSK, QPSK &
Unknown (ISI; composite; blind; partial-posterior); control: Rayleigh &
Extension: unknown-channel failure of AF/DF and learned mitigation (H5) \\
\hline

E7 & SISO & QPSK &
Rayleigh (canonical), rate-1/2 convolutional code &
Supplementary: coded block-DF against symbol-wise DF and coded-aware learned relays (Section~\ref{sec:coded-block-df}) \\
\hline

E8 & SISO & BPSK, QPSK &
All nine families above, plus the coded scenario &
Minimum relay size and the capacity hypothesis H3, with three initializations per configuration (Section~\ref{sec:minimum-relay-size}) \\

\hline
\end{longtable}

\section{Canonical Relay Comparison (SISO, Rayleigh, QPSK)}\label{sec:siso-bpsk-performance-baseline-relay-comparison}

\textbf{Goal:} Evaluate the classical and AI-based relay strategies on the canonical setup of this thesis.

\textbf{Trials:} - \textbf{Topology:} SISO (both hops) - \textbf{Modulation scheme:} QPSK - \textbf{Channel:} i.i.d.\ Rayleigh fast fading (canonical) - \textbf{Equalizers:} One-tap coherent (per hop); no ISI equalizer needed (memoryless) - \textbf{Reported learned relay architectures (six):} Supervised (MLP, Hybrid), Generative (VAE), Sequence (Transformer, Mamba-S6, Mamba-2 SSD) - \textbf{NN activation:} ReLU hidden, tanh output

\textbf{Conclusion:} On the canonical Rayleigh channel, AI relays match but do not surpass symbol-wise DF: DF achieves the lowest or tied-lowest BER at essentially every SNR point, including the low-SNR regime, while the best neural relays (MLP, Hybrid, VAE) track it within statistical uncertainty and clearly outperform AF. This confirms H1 (AI relays beat AF at low SNR) and H2 (DF is the best-performing evaluated symbol-wise classical relay at $\geq 6$ dB, with zero trainable parameters) on the canonical channel.

\subsection{Results Tables}\label{sec:results-tables}

\begingroup\footnotesize
\begin{longtable}[]{@{}lccccccccc@{}}
\caption[BER on the canonical setup]{BER on the canonical setup: QPSK over i.i.d.\ Rayleigh fast fading, SISO on both hops, uncoded, $10\times10{,}000$ bits per SNR point. ``DF th.'' is the closed-form two-hop composition of Eq.~\ref{eq:df-ber-rayleigh} (Section~\ref{sec:rayleigh-two-hop-df-ber}), included for direct comparison against the measured DF column. Column abbreviations: Transf.\ is the Transformer relay, Mamba2 the Mamba-2 SSD relay}\label{tbl:table2}\tabularnewline
\toprule\noalign{}
SNR (dB) & AF & DF & DF th. & MLP & Hybrid & VAE & Transf. & Mamba-S6 & Mamba2 \\
\midrule\noalign{}
\endfirsthead
\toprule\noalign{}
SNR (dB) & AF & DF & DF th. & MLP & Hybrid & VAE & Transf. & Mamba-S6 & Mamba2 \\
\midrule\noalign{}
\endhead
\bottomrule\noalign{}
\endlastfoot
0 & 0.4076 & 0.3337 & 0.3333 & 0.3349 & \textbf{0.3329} & 0.3359 & 0.4014 & 0.4071 & 0.4032 \\
4 & 0.3129 & \textbf{0.2219} & 0.2216 & 0.2233 & 0.2220 & 0.2254 & 0.2758 & 0.2812 & 0.2772 \\
8 & 0.1852 & \textbf{0.1218} & 0.1203 & 0.1229 & 0.1220 & 0.1235 & 0.1424 & 0.1448 & 0.1416 \\
12 & 0.0822 & 0.0580 & 0.0560 & 0.0586 & \textbf{0.0579} & 0.0589 & 0.0628 & 0.0637 & 0.0621 \\
16 & 0.0309 & \textbf{0.0245} & 0.0239 & 0.0247 & \textbf{0.0245} & 0.0247 & 0.0260 & 0.0259 & 0.0251 \\
20 & 0.0110 & \textbf{0.0097} & 0.0098 & 0.0099 & \textbf{0.0097} & \textbf{0.0097} & 0.0101 & 0.0100 & 0.0098 \\
\end{longtable}
\endgroup

Two features of Table~\ref{tbl:table2} deserve comment. First, symbol-wise DF tracks the closed form (the ``DF th.'' column above): reading the abscissa as $E_s/N_0$ and converting to $E_b/N_0$ for QPSK ($k=2$, so $3.01$~dB lower), the two-hop expression $2P(1-P)$ with $P=\tfrac12\bigl(1-\sqrt{\bar\gamma/(1+\bar\gamma)}\bigr)$ gives agreement to within a few percent relative error at every tabulated point ($1.1\%$ at 20~dB, $2.4\%$ at 16~dB, worst case $3.5\%$ at 12~dB, each recomputable from the four-decimal values as printed), which provides a consistency check on the channel, equalization and hard-decision pipeline against the analysis of Section~\ref{sec:rayleigh-two-hop-df-ber}. Figure~\ref{fig:fig10} plots this closed form alongside the measured curves, together with the single-hop floor $P$ itself -- the BER a genie that recovered hop~1 perfectly would still incur on hop~2 -- which is a reference for this uncoded fixed-rate QPSK comparison; no evaluated curve crosses it.

Second, the relays divide by family rather than by paradigm. The three feedforward relays, the MLP, the Hybrid and the VAE, show no difference from DF detectable at this simulation budget from 4~dB upward, nor from each other throughout. That is not a finding of equivalence: as Section~\ref{sec:statistical-significance-testing} states, failure to reject requires a prespecified practical margin and an equivalence test to become one, and neither was performed. The VAE reaches $0.00972$ at 20~dB, matching DF exactly at the resolution of this budget. The two sequence models are the ones that separate, trailing by roughly $0.07$ in BER at 0~dB before converging by 16~dB. Capacity is not the explanation, since the sequence models are the larger of the two groups; the memoryless channel simply offers no temporal structure for their inductive bias to exploit, a reading the equal-parameter comparison of Section~\ref{sec:parameter-normalization-complexity-trade-off} tests directly and confirms.

Two comments on Table~\ref{tbl:table2}. First, symbol-wise DF is the best or tied-best relay at every tabulated point, and the learned relays sit within a few thousandths of it from 4~dB upward: on a matched, known channel the learned relay matches DF rather than beating it, at 169 parameters against DF's none. This is H2, and it is the result the rest of the thesis is measured against. Second, the two sequence models trail the feedforward relays at low SNR (0~dB row of Table~\ref{tbl:table2}) and converge with them by 16~dB, so their extra capacity buys nothing on a memoryless channel; Section~\ref{sec:parameter-normalization-complexity-trade-off} separates that from the parameter-count confound.

\paragraph{Why QPSK and not BPSK.} QPSK places one bit on each axis the Rayleigh coefficient acts in, carries two bits per symbol, and is processed by the I/Q-split relays without modification, which makes it the constellation that exercises the complex-signal processing this thesis studies. Coherent BPSK over the same complex channel would be entirely valid (after channel-phase removal the whole BPSK component lies on the real decision axis, so projecting onto it discards no signal energy) and the choice here is one of experimental design rather than of correctness (Section~\ref{sec:evaluated-configurations}). One consequence must be kept in view when reading these BER figures against BPSK closed forms: the SNR axis here is $E_s/N_0$, and QPSK carries two bits per symbol, so a given abscissa corresponds to an $E_b/N_0$ that is $3.01$~dB lower than the same number on a BPSK axis. The constellations are not on a common axis and their error rates are not directly comparable.

\subsection{Results Figures}\label{sec:results-figures-1}

\begin{figure}[htbp]
\centering
\pandocbounded{\includegraphics[width=\linewidth,height=0.8\textheight,keepaspectratio]{results/canonical_qpsk_rayleigh_ci.png}}
\caption[The canonical setup]{The canonical setup: QPSK over i.i.d.\ Rayleigh fast fading, BER vs.\ $E_s/N_0$ for the eight relays of Table~\ref{tbl:table2}, with 95\% confidence intervals, plus two theoretical reference curves: ``DF (theory)'' (dashed), the closed-form two-hop composition, which the measured DF curve tracks closely, within a few percent relative error everywhere (worst case $3.6\%$, at 12~dB); and the ``single-hop floor'' (dotted), the BER a genie that recovered hop~1 perfectly would still incur on hop~2 -- a reference floor for this fixed-rate, fixed-power relay chain. AF is uniformly worst; no difference between symbol-wise DF and the feedforward learned relays is detected from 4~dB upward at this budget; the two sequence models trail at low SNR and converge by 16~dB}\label{fig:fig10}
\end{figure}

\section{Parameter Normalization \& Complexity Trade-off}\label{sec:parameter-normalization-complexity-trade-off}

This study tests H3 (complexity versus BER) and H4 (architecture convergence at an equal parameter budget) on the canonical setup.

\textbf{Goal:} Isolate architectural inductive biases from parameter count effects and characterize the complexity-performance trade-off for relay denoising.

\textbf{Trials:} - \textbf{Topology:} SISO (both hops) - \textbf{Modulation scheme:} QPSK - \textbf{Channel:} i.i.d.\ Rayleigh fast fading (canonical) - \textbf{Equalizers:} One-tap coherent (per hop); no ISI equalizer needed (memoryless) - \textbf{NN architecture:} All normalized to \textasciitilde3,000 parameters, plus original sizes (169 to 26K) - \textbf{NN activation:} ReLU hidden, tanh output

\textbf{Conclusion.} The 169-parameter MLP is competitive with the larger models in the original canonical comparison. At approximately 3,000 parameters, the cross-model relative BER spread falls from $17.4\%$ at 0~dB to $4.1\%$ at 20~dB. This does not isolate parameter count as the cause: width, depth, state dimensions, input windows, and optimization also differ. Later seed replication shows that the apparent high-SNR convergence is not robust to every initialization (Section~\ref{sec:conclusions}).

\subsection{Results Tables}\label{sec:results-tables-2}

\begingroup\footnotesize
\begin{longtable}[]{@{}lcccccccc@{}}
\caption[Equal-parameter-budget comparison at 3K]{Equal-parameter-budget comparison on the canonical setup (QPSK over Rayleigh). All learned columns are the $\sim$3{,}000-parameter variants; every learned model is scaled to approximately 3{,}000 parameters, so parameter count is removed as an explanation for the residual differences; width, depth, state dimension and window size still co-vary with architecture (Table~\ref{tbl:table28}). AF and DF are the unparameterized references. Column abbreviations: Transf.\ is the Transformer relay, Mamba2 the Mamba-2 SSD relay}\label{tbl:table8}\tabularnewline
\toprule\noalign{}
SNR (dB) & MLP & Hybrid & VAE & Transf. & Mamba-S6 & Mamba2 & AF & DF \\
\midrule\noalign{}
\endfirsthead
\toprule\noalign{}
SNR (dB) & MLP & Hybrid & VAE & Transf. & Mamba-S6 & Mamba2 & AF & DF \\
\midrule\noalign{}
\endhead
\bottomrule\noalign{}
\endlastfoot
0 & 0.3361 & 0.3391 & 0.3424 & 0.4068 & 0.4007 & 0.4001 & 0.4076 & \textbf{0.3337} \\
4 & 0.2244 & 0.2271 & 0.2291 & 0.2818 & 0.2751 & 0.2742 & 0.3129 & \textbf{0.2219} \\
8 & 0.1229 & 0.1229 & 0.1267 & 0.1458 & 0.1406 & 0.1405 & 0.1852 & \textbf{0.1218} \\
12 & 0.0586 & 0.0580 & 0.0604 & 0.0640 & 0.0620 & 0.0622 & 0.0822 & \textbf{0.0580} \\
16 & 0.0247 & \textbf{0.0245} & 0.0251 & 0.0259 & 0.0250 & 0.0251 & 0.0309 & \textbf{0.0245} \\
20 & 0.0099 & \textbf{0.0097} & 0.0101 & 0.0100 & 0.0097 & 0.0098 & 0.0110 & \textbf{0.0097} \\
\end{longtable}
\endgroup

\begin{figure}[htbp]
\pandocbounded{\includegraphics[width=\linewidth,height=0.8\textheight,keepaspectratio]{results/normalized_3k_rayleigh_ber.png}}
\caption[Equal-parameter-budget BER (3K), QPSK/Rayleigh]{BER vs $E_s/N_0$ for all relay variants normalised to $\approx$3{,}000 parameters, QPSK over Rayleigh fading. AF and DF (dashed) are the unparameterised references. Shaded bands show 95\% confidence intervals for the learned relays. The feedforward group (MLP-3K, Hybrid-3K, VAE-3K) tracks DF across the full range; the sequence models (Transformer-3K, Mamba-3K, Mamba2-3K) converge with them only at high SNR}\label{fig:fig-3k-ber}\end{figure}

\section{Minimum Relay Size Across Channel Families}\label{sec:minimum-relay-size}

The complexity study of Section~\ref{sec:parameter-normalization-complexity-trade-off} establishes that no evaluated model improves on the 169-parameter MLP, but leaves two questions open: how far \emph{below} 169 parameters the relay can be taken, and whether the answer is a property of the architecture or of the channel. This section answers both by sweeping relay size downward on every channel family used in this thesis.

\textbf{Goal:} Locate the smallest relay that costs nothing measurable against the thesis's own published relay, per channel, and determine which architectural axis --- window, depth or width --- sets that floor.

\textbf{Trials:} - \textbf{Topology:} SISO, two hops, hop~2 per each channel family's own convention (AWGN or coherently-compensated Rayleigh) - \textbf{Modulation scheme:} BPSK or QPSK, per channel, as used elsewhere in this thesis - \textbf{Channels:} the nine families of Chapters~\ref{sec:experiments} and~\ref{sec:unknown-channels}, plus the coded scenario of Section~\ref{sec:coded-block-df} - \textbf{Comparators:} symbol-wise DF where the channel is memoryless, Viterbi/MLSE where it has memory, block DF for the coded scenario; each channel additionally compared against its own published relay (MLP-169, or MLP-756 coded) - \textbf{Relay architecture:} window $\in\{1,3,5,7\}$, depth $\in\{1,2,3\}$, width from 1 to 48 - \textbf{Replication:} three independent initializations per configuration, each configuration held to its worst - \textbf{Metric:} additional SNR required to reach a given BER

\textbf{Conclusion.} Within the tested grid and degradation criteria, observation-window length is an important resource for the selected three-tap channels. Three of the five memoryless cases meet the zero-degradation criterion with four parameters, whereas the three-tap cases use 73--145. Wider windows improve the reported ISI metric by $8$--$11$~dB in this sweep; added depth does not improve the best tested configuration. These are smallest \emph{tested} configurations under finite training and evaluation budgets, not mathematical lower bounds on architecture size. The supporting table and seed spreads are in Appendix~\ref{sec:appendix-f-minimum-size}.

The full budget-by-budget table is Table~\ref{tbl:table-minsize} in Appendix~\ref{sec:appendix-f-minimum-size}, together with the two figures that show the window effect directly (Figures~\ref{fig:minsize-crossover} and~\ref{fig:minsize-budget}).

\section{Coded Block-DF: Measuring the Remark~\ref{rem:df-terminology} Caveat}\label{sec:coded-block-df}

\paragraph{A note on constellation scope.} The fixed-MCS relay-decoding comparisons here are QPSK throughout. 16-QAM appears only as a rung of the modulation-and-coding ladder in Section~\ref{sec:coded-rate-adaptation}, where the object of study is the rate-and-modulation choice rather than the relay's ability to process a denser constellation; that question is left to future work (Section~\ref{sec:modulation-schemes}).

Remark~\ref{rem:df-terminology} cautions that the symbol-wise DF results ``should not be read as bounds on coded block-DF performance''. That is asserted there and measured here, by a practical finite-length block-DF relay built from a rate-1/2 convolutional code and a soft-decision Viterbi decoder, as a supplementary study on the canonical channel rather than a revision of the core comparison.

\textbf{Setup.} A standard rate-1/2 convolutional code (constraint length $K$, $2^{K-1}$ trellis states, zero-tail terminated per 200-information-bit frame) is applied before QPSK modulation on the canonical Rayleigh channel; \texttt{CodedDecodeAndForwardRelay} decodes the full frame at the relay via soft-decision Viterbi using an unweighted squared-distance metric on the coherently compensated I/Q samples, then re-encodes and re-modulates a fresh codeword before forwarding. Because post-equalization noise is heteroskedastic, a matched metric would weight each symbol by its instantaneous channel reliability. The implemented decoder is therefore a practical mismatched block-DF baseline, not a matched-ML decoder for the fading observations. Two coded-aware learned relays are trained on the same task for comparison: a windowed 756-parameter MLP classifier (window 21, $\approx$5$\times K$ for $K=3$) and the Mamba-S6 architecture already used in Table~\ref{tbl:table2} (24,084 parameters at this configuration), both predicting the transmitted QPSK symbol from a window of noisy coded observations with no explicit knowledge of the code. All results use the thesis-standard budget of 100 trials $\times$ 100,000 information bits per SNR point.

\begin{longtable}[]{@{}lcccccc@{}}
\caption[Coded block-DF vs.\ uncoded DF and coded-aware learned relays]{Coded block-DF vs.\ uncoded symbol-wise DF and coded-aware learned relays, canonical QPSK/Rayleigh, rate-1/2 $K{=}3$ code. ``coded-AF'' forwards the noisy codeword unmodified (the relay does not decode), isolating the code's own gain from any relay intelligence}\label{tbl:table34}\tabularnewline
\toprule\noalign{}
SNR (dB) & uncoded DF & coded AF & coded DF & MLP-coded & Mamba-coded \\
\midrule\noalign{}
\endfirsthead
\toprule\noalign{}
SNR (dB) & uncoded DF & coded AF & coded DF & MLP-coded & Mamba-coded \\
\midrule\noalign{}
\endhead
\bottomrule\noalign{}
\endlastfoot
0 & \textbf{0.3337} & 0.4902 & 0.4204 & 0.4439 & 0.4505 \\
4 & \textbf{0.2216} & 0.4468 & 0.2229 & 0.2974 & 0.3214 \\
8 & 0.1205 & 0.2722 & \textbf{0.0638} & 0.1010 & 0.1174 \\
12 & 0.0560 & 0.0539 & \textbf{0.0154} & 0.0214 & 0.0226 \\
16 & 0.0239 & 0.0099 & 0.0042 & 0.0045 & \textbf{0.0041} \\
20 & 0.0098 & 0.0020 & 0.0014 & 0.0011 & \textbf{0.0010} \\
\end{longtable}

\textbf{Coded-DF is stronger over part of the tested range.} At equal $E_s/N_0$, coded block-DF beats uncoded symbol-wise DF from 8~dB upward, by $5.7\times$ in BER at 16~dB and $7.2\times$ at 20~dB. The rate-1/2 code also reduces information throughput; Section~\ref{sec:coded-latency-throughput} examines an approximately equal-rate comparison. At 0 and 4~dB, coded AF and coded DF are worse than uncoded DF, with the measured crossover between 4 and 8~dB. This threshold is specific to the tested code, metrics, frame and modulation; the experiment does not isolate why it occurs.

\textbf{The larger coded relay does not consistently improve BER.} The convolutional code supplies temporal structure, but both learned relays trail coded-DF at 4--8~dB, and Mamba-S6 has higher measured BER than the 756-parameter MLP there (Table~\ref{tbl:table34}). At higher SNR it can improve BER. Thus this comparison supports neither universal capacity saturation nor a claim that sequence models cannot exploit coding.

\textbf{Paired high-SNR comparisons resolve different orderings.} Table~\ref{tbl:table37} uses 100 paired trials with the same source bits and channel draws for each relay. At 20~dB both learned relays have lower BER than coded-DF in at least 99 trials, with Wilcoxon $p<10^{-4}$. At 16~dB Mamba-coded wins 68 trials, whereas MLP-coded wins 23 and loses 77; both paired comparisons give $p<10^{-4}$. These results concern the implemented configurations, not all possible coded relays.

\begin{longtable}[]{@{}lccccc@{}}
\caption[Paired re-measurement of the high-SNR points]{Paired re-measurement of the two high-SNR points of Table~\ref{tbl:table34}, 100 paired trials $\times$ 100{,}000 information bits. ``Diff'' is the mean per-trial BER difference against coded-DF (negative favours the learned relay); ``W/L'' counts trials won/lost; $p$ is Wilcoxon signed-rank}\label{tbl:table37}\tabularnewline
\toprule\noalign{}
SNR & Relay & BER & Diff vs.\ coded-DF & W/L & $p$ \\
\midrule\noalign{}
\endfirsthead
\toprule\noalign{}
SNR & Relay & BER & Diff vs.\ coded-DF & W/L & $p$ \\
\midrule\noalign{}
\endhead
\bottomrule\noalign{}
\endlastfoot
16 dB & coded-DF & 0.004245 & --- & --- & --- \\
 & MLP-coded & 0.004437 & $+0.000192$ & 23/77 & $<10^{-4}$ \\
 & Mamba-coded & \textbf{0.004124} & $-0.000121$ & 68/32 & $<10^{-4}$ \\
\hline
20 dB & coded-DF & 0.001362 & --- & --- & --- \\
 & MLP-coded & 0.001066 & $-0.000295$ & 99/1 & $<10^{-4}$ \\
 & Mamba-coded & \textbf{0.000967} & $-0.000395$ & 100/0 & $<10^{-4}$ \\
\end{longtable}

\textbf{The observations do not establish a decoding mechanism.} The implemented Viterbi decoder uses unweighted post-equalization distances and is mismatched to symbol-dependent Rayleigh reliability. Table~\ref{tbl:table38} records fewer relay-output symbol errors for coded-DF, but a lower destination repair fraction than for the learned relay. Decode--re-encode forwarding can turn a wrong message into another valid codeword; this is one possible explanation, not a finding that coding redundancy is literally consumed. The experiments do not separate this explanation from metric mismatch, error correlations, code strength, or frame length.

\begin{longtable}[]{@{}lccccc@{}}
\caption[Where the errors go: relay output versus destination output]{Where the errors go, 100 trials $\times$ 500 frames. ``Relay symbol ER'' counts symbols the relay put on the air that differ from the transmitted ones; ``repaired'' is the fraction of those that do not survive the destination's decoder. The oracle relay forwards the clean codeword and so isolates the hop-2-only floor}\label{tbl:table38}\tabularnewline
\toprule\noalign{}
SNR & Relay & Relay symbol ER & Repaired & Final BER \\
\midrule\noalign{}
\endfirsthead
\toprule\noalign{}
SNR & Relay & Relay symbol ER & Repaired & Final BER \\
\midrule\noalign{}
\endhead
\bottomrule\noalign{}
\endlastfoot
16 dB & coded-DF & \textbf{0.00505} & 15.7\% & 0.004257 \\
 & MLP-coded & 0.01065 & \textbf{57.9\%} & 0.004487 \\
 & oracle & --- & --- & 0.002140 \\
\hline
20 dB & coded-DF & \textbf{0.00182} & 24.7\% & 0.001369 \\
 & MLP-coded & 0.00385 & \textbf{73.0\%} & 0.001039 \\
 & oracle & --- & --- & 0.000668 \\
\end{longtable}

\paragraph{Constraint-length sweep.} For the tested $K\in\{3,5,7\}$ generators, larger-$K$ codes have higher BER across 0--8~dB. At 20~dB the rounded means are $0.0014$, $0.0015$, and $0.0012$. The sweep in \url{results/coded_df_experiment.json} does not establish equivalence at 20~dB, nor rule out benefits from other codes, reliability-weighted metrics, or longer frames.

\paragraph{Higher-SNR evaluation.} These short-code measurements do not approach the strong, ideally capacity-achieving code in Remark~\ref{rem:df-terminology}: at 20~dB destination frame errors occur in $20.23\%$ of coded-DF frames and $10.58\%$ even with an oracle relay. Table~\ref{tbl:table44} extends the same code to higher SNR. A second learned relay is retrained on $5$--$30$~dB; its 32-dB point is outside that training range. This varies operating SNR, not code optimality.

\begingroup\footnotesize
\begin{longtable}[]{@{}lccccc@{}}
\caption[Block-DF into the reliable-decoding regime]{Block-DF into the reliable-decoding regime: the rate-1/2 $K{=}3$ code and canonical QPSK/Rayleigh link of Table~\ref{tbl:table34}, extended up the SNR axis, $100$ trials $\times$ $1{,}000$ frames ($20{,}000{,}000$ information bits) per point. ``MLP (5/10/15)'' uses that table's training recipe; ``MLP (5--30)'' is the same architecture retrained on $5$--$30$~dB, bracketing all but the $32$~dB point, hence a different instance whose $20$~dB entry need not reproduce it. Bold marks the best \emph{relay} in each row (both marked where they tie); the oracle forwards the clean codeword and is the single-decode floor they are measured against, not a competitor}\label{tbl:table44}\tabularnewline
\toprule\noalign{}
SNR (dB) & coded-DF & MLP (5/10/15) & MLP (5--30) & oracle & coded-DF FER \\
\midrule\noalign{}
\endfirsthead
\toprule\noalign{}
SNR (dB) & coded-DF & MLP (5/10/15) & MLP (5--30) & oracle & coded-DF FER \\
\midrule\noalign{}
\endhead
\bottomrule\noalign{}
\endlastfoot
20 & 0.001378 & \textbf{0.001064} & 0.001092 & 0.000703 & 0.1998 \\
24 & 0.000485 & \textbf{0.000308} & 0.000318 & 0.000253 & 0.0804 \\
28 & 0.000187 & \textbf{0.000105} & 0.000105 & 0.000092 & 0.0324 \\
32 & 0.000069 & \textbf{0.000037} & \textbf{0.000037} & 0.000034 & 0.0122 \\
\end{longtable}
\endgroup

\textbf{A gap remains in the tested higher-SNR range.} Across 20--32~dB coded-DF's frame-error rate falls from $0.1998$ to $0.0122$, while its BER-to-oracle ratios are $1.96$, $1.92$, $2.03$, and $2.00$. The retrained MLP's ratios are $1.55$, $1.26$, $1.15$, and $1.08$. At 32~dB its 636 frame errors versus the oracle's 607 in 100{,}000 frames do not resolve a difference by the reported $z=+0.8$ comparison; this is not proof of equivalence. These findings use unweighted decoder metrics. Reliability-weighted decoding is needed before attributing the gap to a general disadvantage of block-DF.

\subsection{Soft-Decision Relaying with the Implemented Metrics}\label{sec:coded-soft-decision}

Table~\ref{tbl:table39} compares hard read-out with soft constellation averages. At the classical relay, BCJR replaces Viterbi and supplies symbol weights under the implemented metric; those weights need not be calibrated under fading. For the MLP, the soft average $\sum_k p_k s_k$ replaces the hard argmax without retraining, but requires the weighted-sum computation. Its hard/soft pair shares the same weights, isolating the implemented read-out change rather than proving that the scores are exact posteriors. The paired experiment also includes a clean-codeword oracle relay.

\begin{longtable}[]{@{}lccccc@{}}
\caption[Soft- versus hard-decision relaying on the coded link]{Soft- versus hard-decision relaying, canonical QPSK/Rayleigh, rate-1/2 $K{=}3$, 100 paired trials $\times$ 100{,}000 information bits. The MLP hard/soft pair shares one set of trained weights and differs only in the read-out. The oracle relay forwards the clean codeword, giving the hop-2-only floor. Both block-DF variants are marked at 8~dB, where no difference is resolved}\label{tbl:table39}\tabularnewline
\toprule\noalign{}
SNR (dB) & hard block-DF & soft block-DF & MLP hard & MLP soft & oracle \\
\midrule\noalign{}
\endfirsthead
\toprule\noalign{}
SNR (dB) & hard block-DF & soft block-DF & MLP hard & MLP soft & oracle \\
\midrule\noalign{}
\endhead
\bottomrule\noalign{}
\endlastfoot
0 & 0.4208 & \textbf{0.4150} & 0.4443 & 0.4343 & 0.3022 \\
4 & 0.2231 & \textbf{0.2218} & 0.2969 & 0.2761 & 0.1278 \\
8 & \textbf{0.0636} & \textbf{0.0636} & 0.1008 & 0.0878 & 0.0327 \\
12 & \textbf{0.01537} & 0.01537 & 0.02142 & 0.01768 & 0.00768 \\
16 & 0.004222 & 0.004222 & 0.004403 & \textbf{0.003633} & 0.002105 \\
20 & 0.001354 & 0.002348$^{\dagger}$ & 0.001056 & \textbf{0.000912} & 0.000682 \\
\multicolumn{6}{p{0.94\linewidth}}{\footnotesize $^{\dagger}$Nominal-noise-variance run; miscalibrated at high SNR (see text below). The variance-corrected figure is $0.001376$, not resolved from hard block-DF's $0.001375$.} \\
\end{longtable}

\begin{figure}[htbp]
\centering
\pandocbounded{\includegraphics[width=\linewidth]{results/coded_soft_decision.png}}
\caption[Soft- versus hard-decision relaying: BER curves]{Soft- and hard-decision coded relaying under the implemented metrics, with 100 paired trials. The MLP soft constellation average has lower measured BER than its hard read-out at every sampled SNR. The block-DF curves use the variance-sensitivity correction discussed in the text. Neither the curves nor the repair counts isolate a causal explanation for the high-SNR ordering. The oracle forwards the clean codeword}\label{fig:fig57}
\end{figure}

\textbf{Soft MLP read-out improves the sampled points} (Figure~\ref{fig:fig57}). Its BER is lower than the hard read-out at every measured SNR on identical weights. Here the term posterior mean denotes the softmax-weighted constellation average; calibration was not tested. At 16~dB the hard MLP loses to block-DF in 80 paired trials, wins 19 and ties one, whereas the soft MLP wins all 100. At 20~dB the soft MLP's measured BER is within $0.00023$ of the clean-codeword oracle. The comparison does not identify the physical mechanism or establish optimality of this read-out for destination BER.

\textbf{Soft block-DF does not rescue the implemented block-DF baseline.} The relay is given a nominal noise variance, although post-equalization Rayleigh noise has variance $1/(2\gamma|h|^2)$ per real axis. Its BCJR metrics are therefore mismatched and overconfident on deep-fade symbols. Uniformly inflating the assumed variance removes the 20-dB gap, but this sensitivity sweep does not supply the required per-symbol reliability weights and should not be described as matched calibration. The same limitation applies to unweighted Viterbi distances: multiplying every branch metric by one constant leaves its decisions unchanged, whereas omitted symbol-dependent weights can change the selected path.

Soft and hard block-DF have similar reported BER at 16 and 20~dB under these decoder metrics. This does not isolate hard quantization as the cause of the gap, nor prove a general disadvantage of decoding at a relay.

The coded study therefore supports a configuration-specific high-SNR benefit for the implemented learned relays and a benefit from soft MLP read-out. Since the classical metrics omit per-symbol reliability, it does not establish that denoising at the relay and decoding only at the destination is generally preferable to decode--re-encode forwarding.

\subsection{Throughput and Latency: What the BER Tables Leave Out}\label{sec:coded-latency-throughput}

Every table above compares relays at equal $E_s/N_0$. That is the correct axis for asking which relay detects better, and the wrong one for asking which relay a link should use, because it silently ignores two costs that differ by orders of magnitude across these designs.

\textbf{Throughput: the coded comparison was not like-for-like.} A rate-1/2 code carries half the information per channel use. Concretely, the coded QPSK link spends $202$ symbols to deliver $200$ information bits ($0.99$ bits/symbol) where uncoded QPSK spends $100$ ($2.00$ bits/symbol). The ``$5.7\times$ better at 16~dB'' quoted above is therefore bought with half the data rate, and a link designer given that trade would not read it as a $5.7\times$ improvement at all. The honest comparison holds spectral efficiency fixed, and the measurements already collected supply it: rate-1/2 \emph{16-QAM} delivers $4 \times \tfrac12 \times \tfrac{200}{202} \approx 1.98$ information bits per symbol, matching uncoded QPSK's $2.00$. Table~\ref{tbl:table40} places those two side by side.

\begin{longtable}[]{@{}lccc@{}}
\caption[Coded versus uncoded at equal spectral efficiency]{Coded versus uncoded at equal spectral efficiency ($\approx$2 information bits per channel symbol): uncoded QPSK against rate-1/2 $K{=}3$ 16-QAM, both on the canonical Rayleigh channel. Recomputed from the data of Table~\ref{tbl:table34} and \url{results/coded_df_experiment.json}; no additional simulation}\label{tbl:table40}\tabularnewline
\toprule\noalign{}
SNR (dB) & uncoded QPSK & rate-1/2 16-QAM & Coding gain \\
\midrule\noalign{}
\endfirsthead
\toprule\noalign{}
SNR (dB) & uncoded QPSK & rate-1/2 16-QAM & Coding gain \\
\midrule\noalign{}
\endhead
\bottomrule\noalign{}
\endlastfoot
0 & \textbf{0.3337} & 0.4680 & $0.71\times$ (worse) \\
4 & \textbf{0.2216} & 0.4203 & $0.53\times$ (worse) \\
8 & \textbf{0.1205} & 0.2726 & $0.44\times$ (worse) \\
12 & \textbf{0.0560} & 0.1001 & $0.56\times$ (worse) \\
16 & \textbf{0.0239} & 0.0273 & $0.88\times$ (worse) \\
20 & 0.0098 & \textbf{0.0076} & $1.29\times$ (better) \\
\end{longtable}

The result is a substantial qualification of Section~\ref{sec:coded-block-df}. Held to equal throughput, the rate-1/2 code does not pay for itself anywhere below 20~dB, and where it finally does the gain is $1.29\times$, not the $5.7$--$7.2\times$ the equal-$E_s/N_0$ reading suggested. Coding is not free reliability: it is reliability bought with bandwidth, and on this channel, at this rate and frame length, the purchase is a poor one until the link is already operating well. This does not overturn the earlier section -- coded block-DF genuinely is the stronger \emph{detector}, and Remark~\ref{rem:df-terminology}'s caveat genuinely is real -- but it does bound how much the caveat matters in practice.

\textbf{Relay readiness: block relays buffer substantially more.} The relays differ in how much they must accumulate before producing their first output. Symbol-wise AF and DF can process a current symbol; the symmetric 21-sample learned relay needs $w=10$ future symbols; and block-DF must receive a 202-symbol frame before decoding. These are relay-processing readiness delays, not complete end-to-end latency: in the half-duplex protocol every strategy must still spend channel uses on both source--relay and relay--destination transmissions.

\begin{longtable}[]{@{}lcccc@{}}
\caption[Relay readiness and compute cost of the coded relays]{Relay readiness and measured compute cost. ``Buffer'' is the number of symbols the relay must accumulate before it can produce its first output; it excludes the two half-duplex transmission phases and is not end-to-end latency. Compute is the median of 7 runs over $49{,}894$ symbols after a discarded warm-up, single-threaded CPU. The identical benchmark was run on two different machines, both reported here: neither is more correct than the other, and the pair shows that absolute figures and ratios involving sub-microsecond MLP timings are machine-dependent (\url{results/coded_latency_compute_machines.json})}\label{tbl:table41}\tabularnewline
\toprule\noalign{}
 &  & \multicolumn{2}{c}{$\mu$s/symbol} &  \\
Relay & Buffer (symbols) & Machine A & Machine B & Relative to Viterbi \\
\midrule\noalign{}
\endfirsthead
\toprule\noalign{}
 &  & \multicolumn{2}{c}{$\mu$s/symbol} &  \\
Relay & Buffer (symbols) & Machine A & Machine B & Relative to Viterbi \\
\midrule\noalign{}
\endhead
\bottomrule\noalign{}
\endlastfoot
AF / symbol-wise DF & \textbf{0} & --- & --- & --- \\
MLP hard (756p) & 10 & 0.30 & 0.34 & $50$--$63\times$ faster \\
MLP soft (756p) & 10 & 0.22 & 0.35 & $61$--$68\times$ faster \\
hard block-DF (Viterbi) & 202 & 15.07 & 21.23 & $1\times$ \\
soft block-DF (BCJR) & 202 & 29.30 & 41.21 & $1.94\times$ slower \\
\end{longtable}

The two-machine timings are implementation measurements. BCJR takes approximately $1.94$ times the Viterbi time on each machine ($1.944$ and $1.941$, not equal to three decimals). Absolute decoder times differ by about $40\%$. The MLP timings are too small at this resolution for a precise speedup claim; operation counts and algorithmic readiness are reported separately.

Table~\ref{tbl:table40} limits the coding benefit at approximately equal spectral efficiency for this grid, while Table~\ref{tbl:table41} documents different buffering requirements and reference-program timings. These are separate comparisons: measured milliseconds are not arithmetic counts, and neither yields a general hardware or deployment conclusion.

\subsection{Adaptive Modulation and Coding: Choosing the Rate, Not Just the Relay}\label{sec:coded-rate-adaptation}

Every comparison above, including the equal-throughput one, holds the modulation-and-coding scheme (MCS) fixed and asks which relay achieves the lower BER. A real link does not hold the MCS fixed; it adapts. The question a link actually faces is not ``which relay wins at rate~$R$'' but ``what is the highest \emph{useful} rate this link can sustain at this SNR'', and answering it requires more than one code rate to choose from. The mother code of Section~\ref{sec:coded-block-df} is extended here with the standard punctured rates $2/3$ and $3/4$ (deleting a periodic subset of the rate-1/2 output, re-inserted as soft zeros at the decoder, so the same Viterbi decoder runs unmodified at every rate), and evaluated via bit-interleaved coded modulation across the full grid $\{\text{QPSK}, \text{16-QAM}\} \times \{\text{uncoded}, 1/2, 2/3, 3/4\}$.

The objective is an idealized one-shot \emph{goodput}, $G = R\,(1-\mathrm{FER})$, assuming perfect frame-error detection and discarding failed frames without accounting for feedback or retransmission overhead. $G$ and $R$ are both measured per channel symbol transmitted on the active hop, matching the ergodic-capacity reference $C$ introduced below; because the relay is half-duplex, expressing goodput instead as a fraction of the total two-slot cycle (source-to-relay plus relay-to-destination) would halve every figure in this section without changing which MCS or relay strategy is preferred, since that factor of $\tfrac12$ applies identically to both. Maximizing $G$ over the MCS grid at each SNR traces the link-adaptation envelope, compared here for two implemented classical forwarding strategies: \emph{block-DF}, which decodes the code at the relay and re-encodes its message decision before transmitting, against \emph{denoise-only}, which hard-decides each bit and re-modulates without ever invoking the code, retaining the original transmitted code constraints for destination decoding and requiring no training and no code knowledge at the relay.

\begingroup\footnotesize
\begin{longtable}[]{@{}lccccc@{}}
\caption[Link-adaptation envelope: best MCS per SNR]{Link-adaptation envelope: the MCS maximizing goodput at each SNR, block-DF versus denoise-only relaying, canonical Rayleigh channel, 100 trials $\times$ 200 frames $\times$ 200 information bits per point. Columns are SNR in dB, the goodput-maximizing MCS and its goodput $G$ for each relay, and $C$, the ergodic Rayleigh (Shannon) capacity at that SNR, $C=\log_2(e)\,e^{1/\gamma}E_1(1/\gamma)$ -- an unconstrained-input upper bound no achievable goodput may exceed, included as a sanity ceiling rather than a target: the gap to it here is dominated by the finite constellation and coding loss, not by relay strategy. The capacity column is a \emph{single-link} reference evaluated per active-hop channel symbol; it is not a capacity bound for the half-duplex two-hop relay system studied here, whose own bound would additionally account for the two-slot cycle and the relay constraint}\label{tbl:table42}\tabularnewline
\toprule\noalign{}
SNR & MCS (block-DF) & $G$ & MCS (denoise) & $G$ & $C$ \\
\midrule\noalign{}
\endfirsthead
\toprule\noalign{}
SNR & MCS (block-DF) & $G$ & MCS (denoise) & $G$ & $C$ \\
\midrule\noalign{}
\endhead
\bottomrule\noalign{}
\endlastfoot
 8 & QPSK $\frac{1}{2}$ & \textbf{0.0039} & QPSK $\frac{1}{2}$ & 0.0003 & 2.40 \\
12 & QPSK $\frac{1}{2}$ & \textbf{0.1708} & QPSK $\frac{1}{2}$ & 0.1342 & 3.45 \\
16 & QPSK $\frac{1}{2}$ & 0.5406 & QPSK $\frac{1}{2}$ & \textbf{0.6046} & 4.63 \\
20 & QPSK $\frac{3}{4}$ & 0.9545 & QPSK  $\frac{3}{4}$  & \textbf{0.9635} & 5.88 \\
24 & 16-QAM $\frac{2}{3}$ & \textbf{1.5790} & 16-QAM  $\frac{2}{3}$ & 1.4739 & 7.17 \\
28 & 16-QAM $\frac{3}{4}$ & \textbf{2.3003} & 16-QAM $\frac{3}{4}$ & 2.2625 & 8.48 \\
32 & 16-QAM $\frac{3}{4}$ & \textbf{2.6807} & 16-QAM  $\frac{3}{4}$ & 2.6726 & 9.80 \\
36 & 16-QAM uncoded & 3.3920 & 16-QAM uncoded & 3.3920 & 11.13 \\
40 & 16-QAM uncoded & 3.7336 & 16-QAM uncoded & 3.7336 & 12.46 \\
\end{longtable}
\endgroup

\begin{figure}[htbp]
\centering
\pandocbounded{\includegraphics[width=\linewidth]{results/coded_rate_adaptation.png}}
\caption[Link-adaptation envelope]{The envelope of Table~\ref{tbl:table42} (colored) against every individual fixed MCS on the block-DF relay (grey), which is the visual case for adaptation: the envelope tracks the upper hull of the grey curves rather than any single one of them, switching MCS as SNR crosses each rate's operating threshold. The two relay strategies' envelopes are visually almost the same curve}\label{fig:fig58}
\end{figure}

Two observations follow within this finite modulation-and-coding grid.

\textbf{The envelope is a real staircase, and its top step is ``turn the code off''.} Both relay strategies climb QPSK~1/2~$\to$~QPSK~3/4~$\to$~16-QAM~2/3~$\to$~16-QAM~3/4~$\to$~\emph{16-QAM uncoded} as SNR rises from 8 to 40~dB, each switch occurring once the next rate's threshold has cleared. The optimum at high SNR is no code at all: once the channel is good enough, coding redundancy is pure overhead, and the highest-goodput choice is the raw $4$~bits/symbol of 16-QAM with zero coding loss. This is the same lesson as the equal-throughput correction above, now reached by direct optimization over the grid rather than by hand-picking one comparison. Table~\ref{tbl:table42}'s $C$ column confirms the envelope never approaches the information-theoretic ceiling: at 40~dB the achieved $3.73$ sits at $30\%$ of the $12.46$-bit Shannon capacity there, a gap set by the fixed 16-QAM constellation (itself capped at $4$ bits/symbol) and residual FER, not by anything a smarter relay strategy could close.

\textbf{Once the rate adapts, which relay decodes the code barely matters.} The two envelopes in Table~\ref{tbl:table42} are close at every SNR and identical for $\mathrm{SNR}\geq 36$~dB, where both select 16-QAM uncoded and there is no code for the strategies to differ over. Where they do differ, the gap is small and does not consistently favor one side: block-DF leads at 8, 12, 24, 28 and 32~dB, denoise-only leads at 16 and 20~dB, and the largest gap over the whole sweep is $0.105$ information bits/symbol at 24~dB, under $7\%$ of the envelope value there. Compare this to how much the \emph{MCS} choice moves goodput at fixed relay strategy: at 24~dB the spread across the eight MCS options the block-DF envelope searches over is $1.25$, roughly $12\times$ the $0.105$ gap between relay strategies; at 32~dB and above the MCS spread exceeds $1.7$ while the relay-strategy gap is under $0.01$. Rate adaptation, in other words, dominates the relay-decoding decision by more than an order of magnitude across most of this sweep.

Within the evaluated grid, changing the modulation and code rate has a larger goodput effect than switching between these two forwarding implementations. This is not a neural-versus-classical comparison or a general causal claim about decoding location.

\subsection{Goodput Under an Idealized Buffer-Readiness Budget}\label{sec:coded-latency-capacity}

The coded relay of Section~\ref{sec:coded-block-df} decodes a whole frame before forwarding, and every comparison so far has left the frame length free. A link with a deadline cannot: the deadline bounds the blocklength and with it the redundancy the relay has to spend. The achievable rate once blocklength is fixed is a standard question --- the finite-blocklength capacity of Polyanskiy, Poor and Verd\'{u} \cite{PolyanskiyPoorVerdu2010FiniteBlocklength} and the delay-limited capacity of Hanly and Tse \cite{HanlyTse1998DelayLimited} for fading channels.

The following calculation uses an idealized \emph{buffer-readiness budget}, not physical round-trip latency. It counts one destination frame buffer for both strategies and an additional relay frame buffer for block-DF, while assigning only the learned relay's look-ahead to denoise-only relaying. It omits the channel-use duration of the two half-duplex transmission phases, which both strategies incur. Re-deriving the goodput envelope of Table~\ref{tbl:table42} under this accounting illustrates the effect of relay buffering, but does not establish performance under a complete end-to-end deadline.

\begin{longtable}[]{@{}lcccc@{}}
\caption[Link-adaptation envelope under a 150-symbol readiness budget]{Link-adaptation envelope re-derived under the idealized 150-symbol buffer-readiness budget, same layout and data source as Table~\ref{tbl:table42}}\label{tbl:table43}\tabularnewline
\toprule\noalign{}
SNR (dB) & Best MCS (block-DF) & $G$ (block-DF) & Best MCS (denoise) & $G$ (denoise) \\
\midrule\noalign{}
\endfirsthead
\toprule\noalign{}
SNR (dB) & Best MCS (block-DF) & $G$ (block-DF) & Best MCS (denoise) & $G$ (denoise) \\
\midrule\noalign{}
\endhead
\bottomrule\noalign{}
\endlastfoot
8 & QPSK uncoded & 0.0000 & QPSK uncoded & 0.0000 \\
12 & QPSK uncoded & 0.0000 & QPSK 3/4 & \textbf{0.0105} \\
16 & QPSK uncoded & 0.0251 & QPSK 3/4 & \textbf{0.3582} \\
20 & 16-QAM 3/4 & 0.4966 & QPSK 3/4 & \textbf{0.9635} \\
24 & 16-QAM 3/4 & \textbf{1.5357} & 16-QAM 2/3 & 1.4739 \\
28 & 16-QAM 3/4 & \textbf{2.3003} & 16-QAM 3/4 & 2.2625 \\
32 & 16-QAM 3/4 & \textbf{2.6807} & 16-QAM 3/4 & 2.6726 \\
36 & 16-QAM uncoded & 3.3920 & 16-QAM uncoded & 3.3920 \\
40 & 16-QAM uncoded & 3.7336 & 16-QAM uncoded & 3.7336 \\
\end{longtable}

\begin{figure}[htbp]
\centering
\pandocbounded{\includegraphics[width=\linewidth]{results/coded_latency_capacity.png}}
\caption[Goodput vs. buffer-readiness budget]{Best achievable goodput as a function of the idealized buffer-readiness budget $L_{\max}$, at the four SNRs where the two relay strategies differ (16, 20, 24, 28~dB). Each curve is a step function under this accounting. The omitted two-hop transmission time prevents interpreting this axis as end-to-end latency}\label{fig:fig59}
\end{figure}

\textbf{Under this idealized readiness budget, block-DF can be locked out of coding.} At $L_{\max}=150$ symbols (Table~\ref{tbl:table43}), block-DF's cheapest coded option barely clears the budget while its next-cheapest does not. Under the stated accounting, the 16-dB entries differ by approximately $14.3\times$ ($0.3582$ versus $0.0251$ information bits/symbol), and the 20-dB entries by about $1.9\times$. These figures illustrate buffering sensitivity only; they are not end-to-end latency results.

\textbf{Tighten the budget further and the reversal spreads to SNRs where block-DF led unconstrained.} At $L_{\max}=100$ symbols, block-DF has \emph{no} coded option left at any of these four SNRs -- every one of its coded MCS choices exceeds the budget -- so it is pinned to uncoded transmission throughout, while the denoise-only relay still reaches 16-QAM~2/3 or~3/4 at 20, 24 and 28~dB. At 24 and 28~dB, this is a direct reversal: Table~\ref{tbl:table42}'s unconstrained envelope has block-DF \emph{ahead} at both points, but under the 100-symbol budget denoise-only leads by $1.46\times$ and $1.54\times$ respectively, because the MCS that gave block-DF its unconstrained edge is no longer reachable in time.

The read-out comparison of Section~\ref{sec:coded-soft-decision} does not establish a decoding-location mechanism. The finite-grid results show sensitivity to rate selection and idealized buffering. A hard-deadline deployment recommendation requires a slot-level model of both half-duplex transmissions, processing time and reliability; that analysis is not performed here.
